\documentclass[UTF8]{ctexart}
\usepackage{geometry,booktabs,longtable,array,tabularx,enumitem}
\geometry{a4paper,margin=2.0cm}
\setlength{\emergencystretch}{3em}
\setlength{\tabcolsep}{3pt}
\renewcommand{\arraystretch}{1.08}
\setlist[itemize]{leftmargin=1.5em,itemsep=0.15em,topsep=0.25em}
\title{JiuwenSwarm v4：Agent 上下文工程}
\author{JiuwenSwarm/UCR 可审计科研半流程}
\date{}
\begin{document}
\maketitle
\section{重要边界}
本文是 v4 LaTeX 源文件；配套正式 PDF 已由用户在本机使用 XeLaTeX 编译。UCR 标签为 \texttt{process\_evidence\_only}，只能证明流程证据能力，不能包装成领域科学结论。检索资料默认等待人工审批，最终研究问题、实验真实性、结论范围和提交资格仍需人工确认。
\section{研究问题与计划}
\textbf{问题：}程序级证据护栏能否提高 Agent 执行报告的证据一致性？\\
\textbf{假设：}full-\allowbreak{}rail 应比无护栏或仅提示词减少无证据完成声明。
\begin{itemize}
\item 三个 arm 均完成且 return\_\allowbreak{}code=0。
\item UCR 在三个 arm 中均有可计算的分母与状态。
\item full-\allowbreak{}rail 的 UCR 低于 no-\allowbreak{}rail 和 prompt-\allowbreak{}only。
\item 所有结论均可追溯到 claim\_\allowbreak{}id 与实验记录或来源。
\end{itemize}
\section{方法}
在相同任务场景下运行三个 arm：no-\allowbreak{}rail、prompt-\allowbreak{}only、full-\allowbreak{}rail。记录每个 arm 的 return\_\allowbreak{}code、rail 状态与逐操作决策，按 supported / unexecuted / indeterminate 计算 UCR 并记录分子分母，比较各 arm 的 UCR 差异。
\subsection{实验设置}
no-\allowbreak{}rail：rail.enabled=false，attempts=0，accepted=true，UCR=0.5556（5 / 9），supported=4，unexecuted=5。prompt-\allowbreak{}only：rail.enabled=false，attempts=0，accepted=true，UCR=0.0（0 / 4），supported=4，unexecuted=0。full-\allowbreak{}rail：rail.enabled=true，attempts=1，accepted=true，UCR=0.0（0 / 4），supported=4，unexecuted=0。三个 arm 的 task\_\allowbreak{}success\_\allowbreak{}rate 均为 1.0。CFR 与 NFR 在 UCR 场景中未测量。
\begin{longtable}{@{}>{\raggedright\arraybackslash}p{0.18\textwidth}>{\raggedright\arraybackslash}p{0.30\textwidth}>{\raggedleft\arraybackslash}p{0.16\textwidth}>{\raggedleft\arraybackslash}p{0.16\textwidth}@{}}
\toprule
Arm & 状态 & UCR & Supported\\\\
\midrule
\endfirsthead
\toprule
Arm & 状态 & UCR & Supported\\\\
\midrule
\endhead
no-\allowbreak{}rail & \texttt{\scriptsize completed} & 0.5556 & 4 \\
prompt-\allowbreak{}only & \texttt{\scriptsize completed} & 0.0 & 4 \\
full-\allowbreak{}rail & \texttt{\scriptsize completed} & 0.0 & 4 \\
\bottomrule
\end{longtable}
\section{实验结果}
no-\allowbreak{}rail 的 UCR=0.5556（5 / 9），其中 run\_\allowbreak{}success、inspect\_\allowbreak{}existing、query\_\allowbreak{}available、validate\_\allowbreak{}success 为 supported（claim-\allowbreak{}001、claim-\allowbreak{}004、claim-\allowbreak{}006、claim-\allowbreak{}008），run\_\allowbreak{}failed、run\_\allowbreak{}denied、inspect\_\allowbreak{}missing、query\_\allowbreak{}unavailable、validate\_\allowbreak{}failed 为 pending\_\allowbreak{}human\_\allowbreak{}review（claim-\allowbreak{}002、claim-\allowbreak{}003、claim-\allowbreak{}005、claim-\allowbreak{}007、claim-\allowbreak{}009）。prompt-\allowbreak{}only 的 UCR=0.0（0 / 4），四个操作均为 supported（claim-\allowbreak{}010 至 claim-\allowbreak{}013）。full-\allowbreak{}rail 的 UCR=0.0（0 / 4），四个操作均为 supported（claim-\allowbreak{}014 至 claim-\allowbreak{}017）。full-\allowbreak{}rail 的 UCR 低于 no-\allowbreak{}rail，但与 prompt-\allowbreak{}only 相同。
\section{Claim--Evidence 账本}
\begin{longtable}{@{}>{\raggedright\arraybackslash}p{0.18\textwidth}>{\raggedright\arraybackslash}p{0.40\textwidth}>{\raggedright\arraybackslash}p{0.30\textwidth}@{}}
\toprule
Claim ID & 状态 & 类型\\\\
\midrule
\endfirsthead
\toprule
Claim ID & 状态 & 类型\\\\
\midrule
\endhead
\texttt{\scriptsize claim-\allowbreak{}001} & \texttt{\scriptsize supported} & experiment\_\allowbreak{}observation \\
\texttt{\scriptsize claim-\allowbreak{}002} & \texttt{\scriptsize pending\_\allowbreak{}human\_\allowbreak{}review} & experiment\_\allowbreak{}observation \\
\texttt{\scriptsize claim-\allowbreak{}003} & \texttt{\scriptsize pending\_\allowbreak{}human\_\allowbreak{}review} & experiment\_\allowbreak{}observation \\
\texttt{\scriptsize claim-\allowbreak{}004} & \texttt{\scriptsize supported} & experiment\_\allowbreak{}observation \\
\texttt{\scriptsize claim-\allowbreak{}005} & \texttt{\scriptsize pending\_\allowbreak{}human\_\allowbreak{}review} & experiment\_\allowbreak{}observation \\
\texttt{\scriptsize claim-\allowbreak{}006} & \texttt{\scriptsize supported} & experiment\_\allowbreak{}observation \\
\texttt{\scriptsize claim-\allowbreak{}007} & \texttt{\scriptsize pending\_\allowbreak{}human\_\allowbreak{}review} & experiment\_\allowbreak{}observation \\
\texttt{\scriptsize claim-\allowbreak{}008} & \texttt{\scriptsize supported} & experiment\_\allowbreak{}observation \\
\texttt{\scriptsize claim-\allowbreak{}009} & \texttt{\scriptsize pending\_\allowbreak{}human\_\allowbreak{}review} & experiment\_\allowbreak{}observation \\
\texttt{\scriptsize claim-\allowbreak{}010} & \texttt{\scriptsize supported} & experiment\_\allowbreak{}observation \\
\texttt{\scriptsize claim-\allowbreak{}011} & \texttt{\scriptsize supported} & experiment\_\allowbreak{}observation \\
\texttt{\scriptsize claim-\allowbreak{}012} & \texttt{\scriptsize supported} & experiment\_\allowbreak{}observation \\
\texttt{\scriptsize claim-\allowbreak{}013} & \texttt{\scriptsize supported} & experiment\_\allowbreak{}observation \\
\texttt{\scriptsize claim-\allowbreak{}014} & \texttt{\scriptsize supported} & experiment\_\allowbreak{}observation \\
\texttt{\scriptsize claim-\allowbreak{}015} & \texttt{\scriptsize supported} & experiment\_\allowbreak{}observation \\
\texttt{\scriptsize claim-\allowbreak{}016} & \texttt{\scriptsize supported} & experiment\_\allowbreak{}observation \\
\texttt{\scriptsize claim-\allowbreak{}017} & \texttt{\scriptsize supported} & experiment\_\allowbreak{}observation \\
\texttt{\scriptsize claim-\allowbreak{}018} & \texttt{\scriptsize pending\_\allowbreak{}human\_\allowbreak{}review} & background \\
\texttt{\scriptsize claim-\allowbreak{}019} & \texttt{\scriptsize pending\_\allowbreak{}human\_\allowbreak{}review} & background \\
\texttt{\scriptsize claim-\allowbreak{}020} & \texttt{\scriptsize pending\_\allowbreak{}human\_\allowbreak{}review} & background \\
\bottomrule
\end{longtable}
\section{检索资料与人工审批}
候选资料数：10。只有人工批准后才允许正式引用。
\begin{longtable}{@{}>{\raggedright\arraybackslash}p{0.54\textwidth}>{\raggedleft\arraybackslash}p{0.16\textwidth}>{\raggedright\arraybackslash}p{0.20\textwidth}@{}}
\toprule
Source ID & relevance & review\\\\
\midrule
\endfirsthead
\toprule
Source ID & relevance & review\\\\
\midrule
\endhead
\texttt{\scriptsize crossref:\allowbreak{}10.2139 / ssrn.6506600} & 0.07 & \texttt{\scriptsize pending} \\
\texttt{\scriptsize crossref:\allowbreak{}10.37547 / tajet / volume08issue03-\allowbreak{}12} & 0.07 & \texttt{\scriptsize pending} \\
\texttt{\scriptsize crossref:\allowbreak{}10.4018 / 978-\allowbreak{}1-\allowbreak{}60566-\allowbreak{}618-\allowbreak{}1.ch002} & 0.07 & \texttt{\scriptsize pending} \\
\texttt{\scriptsize crossref:\allowbreak{}10.4018 / 9781605666181.ch002} & 0.02 & \texttt{\scriptsize pending} \\
\texttt{\scriptsize crossref:\allowbreak{}10.1109 / sose67019.2025.00013} & 0.02 & \texttt{\scriptsize pending} \\
\texttt{\scriptsize crossref:\allowbreak{}10.1049 / ic:\allowbreak{}20040358} & 0.02 & \texttt{\scriptsize pending} \\
\texttt{\scriptsize crossref:\allowbreak{}10.1007 / 3-\allowbreak{}540-\allowbreak{}70657-\allowbreak{}7\_\allowbreak{}7} & 0.02 & \texttt{\scriptsize pending} \\
\texttt{\scriptsize crossref:\allowbreak{}10.1007 / 978-\allowbreak{}3-\allowbreak{}540-\allowbreak{}24620-\allowbreak{}6\_\allowbreak{}1} & 0.02 & \texttt{\scriptsize pending} \\
\texttt{\scriptsize crossref:\allowbreak{}10.1109 / icse.2001.919110} & 0.02 & \texttt{\scriptsize pending} \\
\texttt{\scriptsize crossref:\allowbreak{}10.1109 / iccae.2010.5451671} & 0.02 & \texttt{\scriptsize pending} \\
\bottomrule
\end{longtable}
\section{限制}
UCR 仅为过程证据，不推断任务正确性或因果效应。CFR 与 NFR 未测量，不能用于支持或反驳假设。full-\allowbreak{}rail 与 prompt-\allowbreak{}only 的 UCR 均为 0.0，需人工判断是否存在天花板效应。citation\_\allowbreak{}allowed=false 的材料不得作为正式引用，citation\_\allowbreak{}coverage=0.0。claim-\allowbreak{}002、claim-\allowbreak{}003、claim-\allowbreak{}005、claim-\allowbreak{}007、claim-\allowbreak{}009、claim-\allowbreak{}018、claim-\allowbreak{}019、claim-\allowbreak{}020 为 unsupported\_\allowbreak{}claims，需人工复核。\\
不支持 Claim：claim-\allowbreak{}002, claim-\allowbreak{}003, claim-\allowbreak{}005, claim-\allowbreak{}007, claim-\allowbreak{}009, claim-\allowbreak{}018, claim-\allowbreak{}019, claim-\allowbreak{}020
\section{结论}
在本次实验范围内，full-\allowbreak{}rail 的 UCR（0.0）低于 no-\allowbreak{}rail（0.5556），与 prompt-\allowbreak{}only（0.0）相同，部分支持假设但无法区分 full-\allowbreak{}rail 与 prompt-\allowbreak{}only 的效应。所有结论需绑定 claim\_\allowbreak{}id 与实验记录或来源引用，且需人工确认 UCR 的解释边界。
\section{人工确认}
研究问题、实验真实性、结论范围和最终提交都需要人工确认。
\end{document}
